%%=============================================================================
%% Methodologie
%%=============================================================================

\chapter{\IfLanguageName{dutch}{Methodologie}{Methodology}}%
\label{ch:methodologie}
Het onderzoek wordt uitgevoerd aan de hand van een iteratief proces waarin de noden van de betrokken onderzoekers en IT-medewerkers centraal staan. In een eerste fase worden gesprekken gevoerd met een onderzoeker en een teamlead binnen het IT-team. Deze gesprekken hebben als doel om inzicht te krijgen in de huidige werkwijze, de problemen die onderzoekers en ontwikkelaars ervaren en de verwachtingen ten aanzien van mogelijke oplossingen. De resultaten van deze gesprekken worden gebruikt om de onderzoeksvragen verder te concretiseren en om te bepalen aan welke voorwaarden een mogelijke oplossing moet voldoen.

Vervolgens wordt een literatuuronderzoek uitgevoerd naar relevante technieken, methodes en bestaande oplossingen die kunnen bijdragen aan het overbruggen van de kloof tussen de wetenschappelijke analyse en de technische implementatie. Hierbij wordt onder andere gekeken naar bestaande toepassingen en werkwijzen die onderzoekers in staat stellen om modellen, parameters en berekeningen beter te begrijpen, te testen of aan te passen zonder rechtstreeks in de onderliggende software te moeten werken. Het resultaat van deze fase is een overzicht van relevante concepten en mogelijke oplossingsrichtingen die aansluiten bij de eerder vastgestelde noden.

Op basis van de verzamelde informatie worden vervolgens verschillende mogelijke oplossingen uitgewerkt. Hierbij worden de inzichten uit de gesprekken en het literatuuronderzoek gecombineerd om concrete artefacten en proof-of-concepts te ontwikkelen. Deze oplossingen worden niet enkel theoretisch beschreven, maar waar mogelijk ook praktisch uitgeprobeerd binnen de context van het project. Op die manier kan worden nagegaan in welke mate de voorgestelde technieken toepasbaar zijn binnen de bestaande software en onderzoeksworkflow.

De ontwikkelde oplossingen worden vervolgens voorgelegd aan de betrokken onderzoekers en IT-medewerkers. Door middel van feedbackgesprekken wordt nagegaan of de oplossingen aansluiten bij hun noden en verwachtingen en welke verbeteringen of aanpassingen noodzakelijk zijn. De ontvangen feedback wordt gebruikt om de oplossingen verder te verfijnen en opnieuw te evalueren. Hierdoor ontstaat een iteratief proces waarbij de resultaten van een eerdere fase telkens als input dienen voor de volgende fase.

Als eindresultaat van het onderzoek wordt een onderbouwde oplossing of verzameling van aanbevelingen opgeleverd die aantoont hoe de kloof tussen de IT-implementatie en de wetenschappelijke onderzoeksworkflow kan worden verkleind. De uiteindelijke oplossing wordt ge\"evalueerd aan de hand van de vooraf ge\"identificeerde noden en criteria uit de gesprekken en het literatuuronderzoek. Op basis hiervan wordt geconcludeerd in welke mate de voorgestelde aanpak geschikt is voor de context van het project en welke mogelijkheden er zijn voor verdere ontwikkeling.
% <under construction>
% %%%%%%%%%%%%%%%%%%%%%%%%%%%%%%%%%%%%%%%%%%%%
% %% afspraken maken voor gespreken         %%
% %% gepserk met samuel, onderzoeks, mentor %%
% %% schrijven + sturen naar onderzoekers   %%
% %%%%%%%%%%%%%%%%%%%%%%%%%%%%%%%%%%%%%%%%%%%%

% % feedback krijgen & verwerken

% gesprekken met onderzoeker en teamlead van it'ers, kijken wat ze willen, wat de voorstellen zijn

% zoeken naar relevante literatuur onderdelen

% werken aan mogelijke zaken die kunnen helpen

% feedback vragen over die zaken











% Binnen een concrete casus bij het Instituut voor Landbouw-, Visserij- en Voedingsonderzoek (ILVO) wordt een toegepast, onderwerpgericht onderzoeksopzet gevolgd. De focus ligt op het analyseren van de huidige samenwerking tussen onderzoekers en softwareontwikkelaars rond een \text{.NET}-gebaseerd rekenmodel, en op het uitwerken van een haalbare oplossingsrichting die deze samenwerking ondersteunt. Het onderzoekstraject is opgedeeld in vijf fasen, waarbij de feitelijke uitvoering organisch evolueerde op basis van voortschrijdend inzicht.

% \section{Eerste fase: Probleemanalyse}
% In de eerste fase werd het probleemdomein in kaart gebracht door middel van een verkennend gesprek met de copromotor en een grondige analyse van de bestaande documentatie. Aanvankelijk werd de problematiek gedefinieerd als een puur documentatieprobleem, maar gaandeweg werd duidelijk dat de kern van het probleem lag in de afstand tussen de wetenschappelijke logica en de technische implementatie in C\#. Door concrete use-cases van de EMAV-applicatie te bestuderen, werd vastgesteld dat onderzoekers beperkt worden door de \"black-box\" aard van de software.

% \section{Tweede fase: Requirement-analyse}
% De tweede fase richtte zich op het identificeren van de noden van zowel onderzoekers als ontwikkelaars. Voor dit onderzoekstraject is ervoor gekozen om de vereisten primair af te leiden uit de expertise van de copromotor, de analyse van eerdere samenwerkingsmomenten en de functionele behoeften van het EMAV-project. Deze aanpak bood een gerichte basis om de technische requirements te definiëren, waarbij de focus lag op transparantie, autonome experimenteerruimte en het voorkomen van documentation drift, met behoud van de stabiliteit van de productieomgeving.

% \section{Derde fase: Literatuurstudie en technologische verkenning}
% In de derde fase werd gezocht naar technologische oplossingen die de kloof tussen code en documentatie kunnen overbruggen. Er werd een samenvatting opgesteld van de "Stand van zaken", waarbij concepten zoals Literate Programming en interactieve notebooks (zoals Jupyter en Quarto) werden onderzocht. Deze fase was cruciaal om te bepalen welke technologieën compatibel zijn met de bestaande \text{.NET}-architectuur van ILVO en hoe deze ingezet kunnen worden om berekeningen inzichtelijk te maken.

% \section{Vierde fase: Ontwerp en selectie van de oplossingsrichting}
% Op basis van de technologische verkenning en een diepgaande analyse van de EMAV-broncode, werd in de vierde fase een specifieke oplossingsrichting gekozen. Er werd besloten om Python te integreren binnen de \text{.NET}-context. Het ontwerp voorzag in een 'switch'-mechanisme waarmee binnen de applicatie gewisseld kan worden tussen de officiële C\#-implementatie en een experimentele Python-versie. Deze hybride aanpak werd geselecteerd omdat het onderzoekers de vertrouwde Python-omgeving biedt zonder de robuustheid van de C\#-rekenmodules te ondermijnen.

% \section{Vijfde fase: Proof-of-concept en evaluatie}
% In de laatste fase werd de voorgestelde oplossing gerealiseerd in een proof-of-concept. Dit prototype diende als validatie voor de technische haalbaarheid van de Python-integratie en de inzet van notebooks als interface. De evaluatie vond plaats door de deelvragen van de bachelorproef stapsgewijs te beantwoorden op basis van de bevindingen uit de PoC. Hierbij werd getoetst in welke mate de voorgestelde werkwijze effectief bijdraagt aan de autonomie van de onderzoeker en de traceerbaarheid van de berekeningen, wat de basis vormde voor de uiteindelijke conclusies.
